\documentclass[12pt]{jlreq}
\usepackage{amsmath, amssymb}

\begin{document}

\newcommand{\C}{\mathbb{C}}
\newcommand{\R}{\mathbb{R}}
\newcommand{\Z}{\mathbb{Z}}
\newcommand{\Tr}{\operatorname{Tr}}
\newcommand{\Impart}{\operatorname{Im}}
\newcommand{\AF}{\operatorname{AF}}
\newcommand{\AX}{\operatorname{AX}}
\newcommand{\GL}{\operatorname{GL}}

\((\Omega,\mathcal F,P)\) を確率空間とする。
\(X_n,\ n=1,2,3,\ldots\) は同じ分布をもつ独立な確率変数の列で，
\[
  P(X_1\ge0)=1
\]
を満たすとする。また，確率変数 \(Y_n,\ n=1,2,3,\ldots\) を
\[
  Y_n=\sum_{k=1}^{n}X_k^k
\]
で定める。

\begin{enumerate}
  \item \(\displaystyle \sum_{k=1}^{\infty}E[X_1^k]<\infty\) ならば，確率変数列
  \(\{Y_n\}_{n=1}^{\infty}\) は \(n\to\infty\) のとき確率収束することを示せ。
  \item 確率変数列 \(\{Y_n\}_{n=1}^{\infty}\) が \(n\to\infty\) のとき確率収束するならば，
  \[
    \sum_{n=1}^{\infty}P\left(X_1\ge 1-\frac1n\right)<\infty
  \]
  となることを示せ。
  \item 確率変数列 \(\{Y_n\}_{n=1}^{\infty}\) が \(n\to\infty\) のとき確率収束するならば，
  \[
    \sum_{k=1}^{\infty}E[X_1^k]<\infty
  \]
  となることを示せ。
\end{enumerate}

\end{document}
